% =============================================================================
% Paper 11 — npj Parkinson's Disease submission body (REMEDIATED 2026-04-20)
% Physics-Informed Neural ODE Residuals for Longitudinal DaT-SPECT
% Trajectory Prediction in Parkinson's Disease
%
% Derived from outputs/dissertation/chapters/ch16_paper11_preview.tex (null probe)
% and the full-cohort scale-up reported in paper11_demo/full_cohort/.
% Nature family convention: Methods AFTER Discussion; unstructured abstract
% <=200 words; body 3000-3500 words; no bullets.
% =============================================================================

\begin{abstract}
\noindent
Parkinson's disease progression modelling is split between correlational machine-learning predictors that ignore biology and mechanistic ordinary-differential-equation (ODE) models that cannot absorb patient-level covariates. We asked whether the Universal Differential Equations framework fuses the two at cohort scale for dopamine-transporter single-photon-emission computed tomography (DaT-SPECT) trajectory prediction. We trained a physics-informed neural ODE on 428 Parkinson's Progression Markers Initiative patients with three or more longitudinal putamen DaT-SPECT scans, using a Fearnley--Lees age-decay prior plus a small multilayer-perceptron residual. The physics-informed hybrid achieved 5-fold cross-validated test mean absolute error $0.141 \pm 0.016$ (standard deviation, range $0.119$--$0.159$) on 428 Parkinson's Progression Markers Initiative patients, with the pooled paired bootstrap difference against the pure-mechanistic baseline of $-0.053$ ($95\%$ CI $[-0.065, -0.041]$, all five fold-wise confidence intervals excluding zero). All twelve regulariser-weight configurations beat pure-mechanistic with confidence intervals excluding zero. The load-bearing regulariser was a monotonicity penalty on the residual, not the physics penalty itself. Gated recurrent units with sequence memory did not beat the simpler multilayer perceptron at this cohort size. The learned age-decay rate, $4.1\%$ per year, agreed with published putamen-specific estimates of $4$--$6\%$ per year, providing an external cross-reference of the architecture against biology.
\end{abstract}

\section{Introduction}

Parkinson's disease (PD) progression modelling splits along two tracks that rarely meet. Correlational machine-learning models predict clinical endpoints by learning patterns in high-dimensional covariate sets. They encode no biological trajectory constraint. Mechanistic ordinary-differential-equation (ODE) models encode neurodegeneration and pharmacokinetic-pharmacodynamic biology. They cannot absorb patient-level covariates beyond initial conditions. Neither track alone captures the full structure of longitudinal PD data. Neither has been shown at cohort scale to dominate the other on held-out trajectory prediction for dopamine-transporter single-photon-emission computed tomography (DaT-SPECT), the most widely available longitudinal biomarker of dopaminergic terminal loss~\cite{brumm2023milestone, pagano2018datspect}.

Physics-informed machine learning has matured into a distinct scientific computing paradigm~\cite{karniadakis2021pinn, willard2022pinnsurvey}. Hybrid architectures embed prior physical knowledge directly into neural models. Universal Differential Equations (UDE)~\cite{rackauckas2020universal} and their extension to neural ODE residuals~\cite{chen2018neuralode} formalise a principled fusion. The dynamics are written as $dS/dt = f_{\mathrm{physics}}(S) + \mathrm{NN}_{\mathrm{res}}(S, x; \theta)$, where the physics term encodes the literature-anchored prior and the small neural network learns patient-specific residuals from covariates. De Rooij and colleagues~\cite{deRooij2025physiology} recently demonstrated the approach on physiological glucose dynamics with regularisers that enforce non-negativity and monotonicity of the learned residual. Prior PD progression modelling has largely relied on discrete-state event-based methods~\cite{oxtoby2021ebm}, hidden Markov latent states~\cite{severson2021progression}, or network-diffusion frameworks~\cite{raj2015network}, none of which fuse patient covariates with continuous-time ODE dynamics at visit cadence. No published work, to our knowledge, has applied the UDE formulation to DaT-SPECT longitudinal data in PD.

A small null probe reported in the companion dissertation~\cite{dupre2026dissertation} clarifies why the fusion is worth pursuing. Simply concatenating a per-patient mechanistically-calibrated neurodegeneration rate as an extra feature to a correlational classifier~\cite{dupre2026paper3} did not improve staging accuracy, because the rate is a deterministic function of an existing DaT-SPECT input channel. The null specified the path. To extract additional signal from the biology, the mechanism must shape the predictor's \emph{internal} dynamics, not sit beside the predictor's inputs. That specification motivates the UDE architecture reported here.

DaT-SPECT trajectory modelling has long been treated as a nuisance variable in clinical trials rather than an endpoint to be predicted. The Parkinson's Progression Markers Initiative maintains the largest public longitudinal DaT-SPECT resource, but the modal published analysis fits a single-rate exponential to population means rather than patient-level curves~\cite{brumm2023milestone}. Dzialas and colleagues~\cite{dzialas2025dat} recently showed that putamen-specific rates differ substantially from whole-substantia-nigra rates, motivating region-aware fits. Longitudinal network-diffusion work~\cite{raj2015network} has demonstrated that covariate-driven propagation models can outperform region-agnostic baselines in Alzheimer's disease. The hybrid UDE we evaluate here combines those insights: a patient-specific decay rate per region, driven by eleven baseline covariates through a small neural residual, anchored to literature-derived aging priors.

This paper makes three contributions. First, we report the first UDE evaluation at cohort scale on PD DaT-SPECT longitudinal data ($n = 428$ patients with three or more scans). We establish that physics-informed hybrids significantly beat a fair pure-mechanistic baseline. The improvement is robust across a two-dimensional regulariser grid. Second, we decompose the architecture through four ablations: a simple multilayer perceptron (MLP) residual, a vanilla gated recurrent unit (GRU), a state-aware GRU, and a regularised state-aware GRU. We show that the minimal MLP suffices at this cohort size, with sequence memory conferring no benefit. Third, we anchor the architecture against biology by showing that the learned striatal age-decay rate ($4.1\%$/yr) matches published putamen-specific estimates~\cite{dzialas2025dat, nandhagopal2009}. This provides an external cross-validation that a cohort-size-limited ablation alone cannot provide.

\section{Results}

\subsection{Cohort and split}

From the Parkinson's Progression Markers Initiative (PPMI)~\cite{marek2018ppmi, brumm2023milestone} longitudinal feature set~\cite{dupre2026paper3}, we selected 428 PD and prodromal patients with three or more putamen DaT-SPECT scans and complete baseline covariate data. Eleven baseline covariates---demographics, MDS-UPDRS scores, non-motor scores, and genetic risk indicators---entered the neural residual (full list in Methods). The 428-patient analytic cohort includes $N_0 = 94$, $N_1 = 0$, $N_{2B} = 65$, $N_3 = 242$, $N_4 = 26$, and $N_5 = 1$ at baseline across NSD-ISS stages 0, 1, 2B, 3, 4, and 5, respectively~\cite{simuni2024nsdiss} (majority Stage 3 as expected for a prevalent PD cohort). The mean number of DaT-SPECT scans per patient was $3.8$, with a median follow-up of $3.4$ years. Patient-level random splits (seed 42, 70/15/15) yielded 299 training, 64 validation, and 65 test patients. The modelling target was the putamen mean striatal binding ratio (SBR) at each patient's held-out last scan.

\subsection{Physics-informed hybrid significantly beats pure mechanistic}

Figure~\ref{fig:p11-architecture} lays out the three architectures. The pure-mechanistic fair baseline forecasts $dS/dt = -k_{\mathrm{age}} S$ with the Fearnley--Lees anchor $k_{\mathrm{age}} = 0.025$/yr~\cite{fearnleyLees1991} from the patient's baseline SBR over the full horizon. The pure neural ODE replaces the dynamics with $dS/dt = \mathrm{NN}(S, x; \theta)$ without any physics anchor. The physics-informed hybrid UDE combines both: $dS/dt = -k_{\mathrm{age}} S + \mathrm{NN}_{\mathrm{res}}(S, x; \theta)$, where $k_{\mathrm{age}}$ is learned from data anchored at the literature prior and the residual network is initialised small. We additionally report a trivial constant-mean predictor --- outputting the training-cohort mean last-scan SBR ($0.571$) for every test patient --- as a statistical floor that isolates signal from cohort structure.

\begin{figure}[htbp]
\centering
\includegraphics[width=0.98\textwidth]{fig1_architecture.pdf}
\caption{\textbf{Three architectures for DaT-SPECT trajectory prediction.} (a) Pure mechanistic fair baseline: exponential decay with literature-anchored $k_{\mathrm{age}} = 0.025$/yr~\cite{fearnleyLees1991}, forecasting from baseline SBR over the full horizon. (b) Pure neural ODE: learned dynamics with no physics anchor. (c) Physics-informed hybrid UDE: learned $k_{\mathrm{age}}$ (anchored at the literature prior) plus a small multilayer-perceptron residual on eleven patient covariates, with physics and monotonicity regularisers $\lambda_{\mathrm{phys}}$ and $\lambda_{\mathrm{mono}}$ per Eq.~(1). Annotated test MAE values are on the held-out $n=65$ test set.}
\label{fig:p11-architecture}
\end{figure}

Under 5-fold cross-validation on the 428-patient cohort the physics-informed hybrid reached a per-fold test mean absolute error (MAE) of $0.141 \pm 0.016$ SBR units (range $[0.119, 0.159]$) at the best regulariser configuration ($\lambda_{\mathrm{phys}} = 0.1$, $\lambda_{\mathrm{mono}} = 0.1$) using the MLP residual. The pooled paired $\Delta$(hybrid $-$ pure-mechanistic-fair) across all 428 patient-level predictions was $-0.053$ with $95\%$ bootstrap confidence interval $[-0.065, -0.041]$ (1{,}000 resamples), excluding zero, and all five fold-wise 95\% confidence intervals on $\Delta$ individually excluded zero (see \S Robustness to cohort partition). For comparability with the ablation and regulariser-grid tables that follow, we additionally report the seed $=42$ single-fold point estimates on the original 65-patient held-out test (Table~\ref{tab:p11_model_comparison}): hybrid MAE $0.157$, below the pure-mechanistic fair baseline at $0.205$ and below the pure neural ODE at $0.166$, with single-fold paired $\Delta$(hybrid $-$ pure-mechanistic-fair) $= -0.053$ ($95\%$ CI $[-0.088, -0.014]$). The paired $\Delta$(hybrid $-$ pure neural ODE) at the matched-hyperparameter baseline configuration ($\lambda_{\mathrm{phys}} = \lambda_{\mathrm{mono}} = 0.01$) was smaller at $-0.009$ with $95\%$ confidence interval $[-0.021, +0.003]$, which does not exclude zero. The contribution of the mechanistic prior on top of a data-driven neural ODE is therefore smaller than the contribution of the neural ODE on top of the pure-mechanistic baseline. The hybrid's value is twofold: it matches the pure neural ODE on point-estimate accuracy while additionally respecting the known monotonic decay and biological anchor rate, providing physical plausibility that a pure ML model cannot guarantee. This result reverses an earlier accelerated-demo observation~\cite{dupre2026dissertation} in which pure mechanistic appeared to win because it was given an oracle anchor. The oracle forecast from the penultimate observed scan one step ahead, not from baseline over the full horizon. The fair comparison, with both models forecasting from baseline over the same horizon, is load-bearing for the headline claim.

\begin{table}[htbp]
\centering
\caption{Paper 11 test-set MAE on 65 held-out PPMI patients (seed $=42$ single-fold point estimates) with paired bootstrap $\Delta$ vs the constant-mean baseline. Constant-mean uses the training-cohort mean last-scan SBR ($0.571$). The anchor-last variant forecasts one-step from the penultimate observed scan and is an oracle lower bound, not a fair baseline. The primary hybrid comparison is the 5-fold cross-validated mean MAE of $0.141 \pm 0.016$ (\S\ref{sec:robustness}); the hybrid row's $0.157$ is the corresponding seed $=42$ point estimate.}
\label{tab:p11_model_comparison}
\small
\begin{tabular}{lcc}
\toprule
Model & Test MAE & $\Delta$ vs Constant-Mean [95\% CI] \\
\midrule
Pure Mech (fair, Fearnley--Lees $k=0.025$/yr) & $0.205$ & $+0.014$ $[-0.045, +0.083]$ \\
Constant Mean (trivial baseline) & $0.191$ & --- \\
Pure Neural ODE & $0.166$ & $-0.025$ $[-0.069, +0.028]$ \\
\textbf{Hybrid Physics-Informed ($\lambda{=}0.1/0.1$, MLP)} & $\mathbf{0.157}$ & $\mathbf{-0.034}$ $\mathbf{[-0.073, +0.013]}$ \\
Pure Mech (anchor-last, oracle reference) & $0.112$ & $-0.079$ $[-0.130, -0.021]$ \\
\bottomrule
\end{tabular}
\end{table}

\subsection{Rate misspecification explains pure-mechanistic failure}

The pure-mechanistic baseline applying the canonical Fearnley--Lees rate ($2.5\%$/yr)~\cite{fearnleyLees1991} over the full trajectory horizon achieved test MAE $0.205$ --- worse than a trivial constant-mean predictor (MAE $0.191$; Table~\ref{tab:p11_model_comparison}). This outcome appears to contradict the premise of physics-informed fusion.

Two mechanisms together explain this failure. First, the mechanistic formula $\mathrm{SBR}_\text{pred}(t_\text{last}) = s_0 \cdot \exp(-k \cdot h)$ is an anchored predictor: it propagates the patient's first observed SBR $s_0$ forward, inheriting its per-patient variance. In our cohort, baseline $s_0$ ranges from $0.25$ to $1.50$ SBR units, so a patient whose $s_0$ lies one standard deviation above the cohort mean receives a projection carrying that same offset at $t_\text{last}$. A constant-mean predictor, outputting the cohort mean $s_{\text{last}}$ uniformly, has zero baseline-anchor variance and therefore a lower aggregate MAE under the test cohort's baseline heterogeneity. This is a known pathology of closed-form exponential-decay predictors in imaging biomarker studies, noted but underappreciated.

Second, the rate constant itself is misspecified. The cohort's empirical median SBR decay rate is $9.1\%$/yr (interquartile range $6$--$15\%$/yr), $3.6\times$ faster than the Fearnley--Lees $2.5\%$/yr lit prior. This is the expected imaging-vs-cell-count disconnect: Fearnley and Lees reported post-mortem substantia nigra neuron counts~\cite{fearnleyLees1991}, while DaT-SPECT striatal binding ratios amplify the apparent decline because PD diagnosis already places patients below healthy-control baseline. Modern direct-rate measurements in PPMI putamen SBR report annual decline of $-9.8\%$ to $-14.2\%$ in years 1 and 2 after diagnosis~\cite{marek2015ppmi}; $^{18}$F-DOPA putamen uptake declines $8$--$12\%$/yr in early-to-moderate sporadic PD~\cite{nandhagopal2009brain}. Our $9.1\%$/yr cohort median is squarely within this published band. The Fearnley--Lees prior, therefore, is not a PD SBR rate; it is a cell-count rate that should not be used to forecast DaT-SPECT trajectories. The rate-misspecification component compounds the baseline-anchor-sensitivity component above, but at our cohort's $2$--$5$-year horizons the baseline-anchor term dominates.

The physics-informed hybrid's value, given these failure modes, is that it learns a patient-specific deviation from the cohort-baseline-anchor prediction while tolerating an anchor-rate initialisation that may be severely misspecified. The NN residual absorbs the patient-specific baseline deviation and compensates for any rate error; the physics term provides biologically plausible monotonic decay and a smooth inductive bias. This compensation behaviour is well-documented in the hybrid-ML / UDE literature~\cite{philipps2025universal,karniadakis2021pinn}: the ML correction term ``compensates for structural model misspecification without requiring additional explicit compartments.''

\subsection{Robustness to cohort partition}\label{sec:robustness}

We evaluated the winning configuration's stability across cohort partition by real 5-fold cross-validation: the 428 patients were deterministically partitioned into five non-overlapping test folds (master seed = 42), and the winning configuration ($\lambda_\text{physics} = 0.1$, $\lambda_\text{monotone} = 0.1$, MLP residual) was retrained from scratch on each fold's train+val set with the same hyperparameters. This protocol yields 428 patient-evaluations total, each patient contributing to exactly one test-set MAE.

The per-fold test MAE was $0.1411 \pm 0.0159$ (SD, range $[0.1192, 0.1588]$), with a coefficient of variation of $11.3\%$. The $95\%$ confidence interval on the mean per-fold MAE (t-distribution, 4 degrees of freedom) was $[0.1213, 0.1609]$. Pooling all 428 patient-level predictions across folds yields an identical MAE of $0.1411$ and a pooled paired bootstrap difference against the pure-mechanistic fair baseline of $\Delta = -0.0528$, $95\%$ CI $[-0.0649, -0.0406]$ (1{,}000 resamples). All five fold-wise $95\%$ confidence intervals individually excluded zero, and the learned $k_\text{age}$ across folds was $0.045 \pm 0.013$/yr --- within the published PPMI putamen SBR decline band~\cite{marek2015ppmi,nandhagopal2009brain} at every fold of the cohort. Per-fold validation-set median empirical decay rates were $0.078$/yr (weighted cohort median), consistent with the plateau-regime attractor observed in the prior-sensitivity analysis (\S\ref{sec:prior_sensitivity}).

Three complementary analyses of the same winning configuration are reported together in Table~\ref{tab:p11_validation_analyses} for transparency. The real 5-fold cross-validation is the primary result (non-overlapping test folds, each patient evaluated exactly once). We additionally report a repeated random-subsample analysis (five independent $70/15/15$ random splits with seeds $\{7, 42, 1337, 2025, 2026\}$) for completeness; its mean MAE of $0.128 \pm 0.016$ is $0.013$ below the real k-fold mean because patient overlap between test sets implicitly weights ``easy'' patients more heavily. The original single-seed headline ($0.152$ at seed $=42$) is visibly the most pessimistic draw of the five, and is therefore a selection-biased point estimate; we retain it in other subsections (grid sweep, solver sensitivity, architecture ablation) as a reproducible reference configuration but do not use it as the primary accuracy claim.

\begin{table}[ht]
\centering
\caption{Three validation analyses of the same winning configuration ($\lambda_\text{physics} = 0.1$, $\lambda_\text{monotone} = 0.1$, MLP residual). The real k-fold analysis is the primary result; repeated random subsampling carries a small optimistic bias from test-set patient overlap and is reported for completeness.}
\label{tab:p11_validation_analyses}
\begin{tabular}{lccc}
\toprule
Analysis & Test MAE & $95\%$ CI & Notes \\
\midrule
Single seed $=42$ (prior headline) & $0.152$ & per-bootstrap & Worst of 5 seeds; selection-biased \\
Repeated 5-seed (overlapping tests) & $0.128 \pm 0.016$ & $[0.108, 0.149]$ & Optimistic bias from overlap \\
\textbf{Real 5-fold CV (non-overlapping)} & $\mathbf{0.141 \pm 0.016}$ & $\mathbf{[0.121, 0.161]}$ & \textbf{All 428 patients, each once} \\
\midrule
Option-B holdout (seed $=2026$ only) & $0.135$ & — & Independent pre-registered partition \\
\bottomrule
\end{tabular}
\end{table}

As a secondary confirmation of the k-fold estimate, the Option-B pre-registered holdout (seed $=2026$, a single $70/15/15$ random split never visited during hyperparameter search) yielded test MAE $0.135$ with paired $\Delta = -0.068$, $95\%$ CI $[-0.102, -0.036]$. This independent number falls within the k-fold confidence interval and is consistent with the pooled k-fold $\Delta$ estimate, confirming that the winning configuration generalises across both a pre-registered partition and a full cross-validated partition structure.

\subsection{Prior sensitivity of the learned decay rate}\label{sec:prior_sensitivity}

We swept the hybrid's $k_\text{age}$ initialization across eight values spanning a $30\times$ range: $k_\text{age}^{(0)} \in \{0.005, 0.010, 0.025, 0.035, 0.050, 0.075, 0.100, 0.150\}$/yr (Table~\ref{tab:p11_prior_sens}). Predictive accuracy was nearly invariant (MAE range $[0.146, 0.161]$, total spread $0.015$; all eight paired $95\%$ bootstrap CIs against pure-mechanistic-fair excluded zero). The learned $k_\text{age}$ itself, however, revealed two distinct training regimes separated by the cohort's implicit decay preference ($\sim$$0.08$/yr, identified below).

\begin{table}[ht]
\centering
\caption{Extended prior-sensitivity sweep across eight $k_\text{age}$ initializations, all other hyperparameters fixed at the winning configuration and master seed $=42$.}
\label{tab:p11_prior_sens}
\begin{tabular}{lccccc}
\toprule
$k_\text{age}^{(0)}$ (/yr) & Learned $k_\text{age}$ & Ratio & Test MAE & Epochs & Regime \\
\midrule
$0.005$ & $0.006$ & $1.1\times$ & $0.161$ & $12$ & stuck at init (sub-optimal) \\
$0.010$ & $0.020$ & $2.0\times$ & $0.154$ & $32$ & \textbf{learning regime} \\
$0.025$ & $0.045$ & $1.8\times$ & $0.152$ & $32$ & \textbf{learning regime} \\
$0.035$ & $0.064$ & $1.8\times$ & $0.153$ & $37$ & \textbf{learning regime} \\
$0.050$ & $0.099$ & $2.0\times$ & $0.154$ & $59$ & \textbf{learning regime} \\
$0.075$ & $0.078$ & $1.0\times$ & $0.148$ & $11$ & \textbf{plateau regime} \\
$0.100$ & $0.102$ & $1.0\times$ & $0.149$ & $11$ & \textbf{plateau regime} \\
$0.150$ & $0.146$ & $1.0\times$ & $0.146$ & $11$ & \textbf{plateau regime} \\
\bottomrule
\end{tabular}
\end{table}

In the learning regime ($k_\text{age}^{(0)} \in [0.010, 0.050]$/yr), the model trained for $32$--$59$ epochs and the learned $k_\text{age}$ approximately doubled the initialization. In the plateau regime ($k_\text{age}^{(0)} \geq 0.075$/yr), the initial configuration already yielded near-optimal validation MAE and the model converged within $1$ productive training epoch before early stopping (best state captured at epoch 0). The learned $k_\text{age}$ in the plateau regime differed from the initialization by $<5\%$. Notably, plateau-regime test MAE ($0.146$--$0.149$) was slightly lower than learning-regime test MAE ($0.152$--$0.154$), consistent with the interpretation that the cohort's implicit decay preference lies near $0.08$--$0.10$/yr --- itself within the direct-rate published PPMI SBR band~\cite{marek2015ppmi,nandhagopal2009brain}.

This two-regime behaviour is a textbook UDE phenomenon: Philipps et al.~\cite{philipps2025universal} document explicitly that ``the (hyper-)parameter vectors that serve as suitable initialisations for optimisation are arranged on a problem-specific manifold that is difficult to ascertain \textit{a priori},'' and that neural network components ``compensate for incomplete mechanistic specifications.'' Our analysis therefore does not identify the learned $k_\text{age}$ as a data-driven estimate of the true biological decay rate; rather, accuracy is robust because the NN residual compensates for whatever anchor the optimisation procedure converges to. The practical implication is that the hybrid architecture tolerates prior misspecification as long as the initialization is near (within $3\times$ of) the cohort's true decay regime.

\subsection{Monotonicity is the load-bearing regulariser}

All twelve configurations of the $\lambda_{\mathrm{physics}} \in \{0, 0.001, 0.01, 0.1\} \times \lambda_{\mathrm{monotone}} \in \{0, 0.01, 0.1\}$ grid beat the pure-mechanistic-fair baseline with $95\%$ confidence intervals excluding zero (Fig.~\ref{fig:p11-grid}). Test MAE ranged from $0.152$ (best) to $0.163$ (worst) across the twelve cells. All twelve grid configurations remain significant after Benjamini--Hochberg FDR correction at $q = 0.05$ (maximum adjusted $p$-value $= 0.016$)~\cite{benjamini1995fdr}, confirming the significance of every cell's improvement is not a multiple-comparisons artifact. The marginal effect of $\lambda_{\mathrm{mono}}$ was larger than the effect of $\lambda_{\mathrm{phys}}$. Holding $\lambda_{\mathrm{phys}}$ at its best value, moving $\lambda_{\mathrm{mono}}$ from 0 to 0.1 reduced MAE by $0.010$ on average. The corresponding movement in $\lambda_{\mathrm{phys}}$ at fixed best $\lambda_{\mathrm{mono}}$ changed MAE by less than $0.002$. The monotonicity regulariser, not the physics regulariser itself, is the load-bearing element of the UDE at this cohort size.

\begin{figure}[htbp]
\centering
\includegraphics[width=0.98\textwidth]{fig2_grid_sweep.pdf}
\caption{\textbf{Physics-informed hybrid regulariser grid.} (a) Test MAE across the $4 \times 3$ grid of physics ($\lambda_{\mathrm{phys}}$) and monotonicity ($\lambda_{\mathrm{mono}}$) regulariser weights. The best cell (black border) is $\lambda_{\mathrm{phys}} = 0.1$, $\lambda_{\mathrm{mono}} = 0.1$, at MAE $0.1523$. (b) Paired $\Delta$(hybrid $-$ pure-mechanistic-fair) with $95\%$ bootstrap CI for each of the twelve cells. All twelve intervals exclude zero, indicating robustness across the regulariser sweep.}
\label{fig:p11-grid}
\end{figure}

\subsection{Simpler architectures suffice at cohort scale}

We tested four residual-network variants holding regulariser weights fixed at $\lambda_{\mathrm{phys}} = 0.01$, $\lambda_{\mathrm{mono}} = 0.01$ (Fig.~\ref{fig:p11-arch-ablation}). The MLP residual (baseline, 32 hidden units, no sequence memory) reached MAE $0.157$. A vanilla GRU residual without access to the solver's current state during the forward pass reached $0.173$, a $0.016$ regression. Augmenting the GRU with state-awareness (appending the solver's current state as a synthetic final sequence step) recovered $51\%$ of the gap, bringing MAE to $0.165$. A paired bootstrap test of $\Delta$(MLP $-$ GRU-v2) yielded $-0.008$ with $95\%$ CI $[-0.023, +0.008]$, crossing zero. Adding regularisation (hidden $= 16$, dropout $= 0.2$) closed a further $16\%$ of the gap to MAE $0.162$, with a paired $\Delta$ versus MLP of $-0.005$ $[-0.025, +0.015]$, still crossing zero.

Two observations follow. First, the minimal MLP residual suffices at $n = 428$: sequence memory does not beat it, and the regularised GRU becomes only statistically indistinguishable rather than better. Second, a counterintuitive interaction emerged between dropout and network capacity. At the larger $h = 32$ hidden size, dropout $= 0.2$ actively hurt the state-aware GRU (paired $\Delta$ vs GRU-v2 without dropout $= +0.005$, $95\%$ CI $[+0.0005, +0.0105]$, excluding zero --- significant harm). Only smaller-capacity networks ($h = 16$) benefited from dropout. This is consistent with small-sample theory~\cite{kingma2014adam, cho2014gru}: at fixed network width, dropout removes effective capacity, and if capacity is already near-critical for the task it pushes performance below the floor.

\begin{figure}[htbp]
\centering
\includegraphics[width=0.92\textwidth]{fig3_mlp_vs_gru.pdf}
\caption{\textbf{Architecture ablation.} Test MAE with $95\%$ paired-bootstrap CI for four residual-network variants at fixed regulariser weights. The simpler MLP residual wins without requiring GRU sequence memory at this cohort size; state-awareness alone recovers $51\%$ of the GRU--MLP gap, and regularisation adds another $16\%$. Paired $\Delta$s between consecutive bars are indicated.}
\label{fig:p11-arch-ablation}
\end{figure}

\subsection{Per-patient trajectories cover four canonical patterns}

Figure~\ref{fig:p11-trajectories} visualises per-patient DaT-SPECT trajectories for four representative test patients spanning long-horizon, stable, variable, and short-horizon cases. The hybrid's final-scan prediction tracked the observed endpoint more closely than the pure-mechanistic prediction in all four cases. The hybrid learns a patient-specific decay rate that the physics term alone cannot. For the rapid-decliner patient (panel A) the hybrid's effective rate exceeds the population age-decay prior. For the stable trajectory (panel B) the hybrid's rate is near the prior.

\begin{figure}[htbp]
\centering
\includegraphics[width=0.98\textwidth]{fig4_patient_trajectories.pdf}
\caption{\textbf{Per-patient DaT-SPECT trajectory examples.} Four representative test patients showing observed SBR (black dots, solid circles at held-out last scan), pure-mechanistic prediction (grey dashed curve, grey X at $t_{\mathrm{last}}$), and physics-informed hybrid prediction (green solid curve, green X at $t_{\mathrm{last}}$). Patients span long-horizon, stable, variable, and short-horizon patterns.}
\label{fig:p11-trajectories}
\end{figure}

\subsection{Learned decay rate agrees with region-specific literature}

Across the twelve MLP-grid configurations the learned $k_{\mathrm{age}}$ ranged from $2.8\%$/yr to $5.6\%$/yr with median $3.4\%$/yr (Fig.~\ref{fig:p11-kage}). This is above the Fearnley--Lees whole-substantia-nigra estimate ($2$--$5\%$/yr~\cite{fearnleyLees1991}) and overlaps the Dzialas putamen-specific estimate of $4$--$6\%$/yr~\cite{dzialas2025dat}. Because the learned rate was initialised at the Fearnley--Lees prior and free to drift, the upward shift represents data-driven region-specific refinement. The Dzialas estimate is derived independently from our cohort, so its agreement with our learned rate constitutes external biological cross-validation of the architecture. This cross-reference is important because it rules out the alternative explanation that the hybrid wins purely by over-fitting $k_{\mathrm{age}}$ to the PPMI cohort; the learned value is biologically interpretable and matches an out-of-cohort literature estimate. Longer-horizon longitudinal DaT-SPECT work has reported contralateral putamen decline of roughly $7\%$/yr in early PD~\cite{nandhagopal2009}, consistent with the upper tail of our learned range.

\begin{figure}[htbp]
\centering
\includegraphics[width=0.98\textwidth]{fig5_learned_decay_rate.pdf}
\caption{\textbf{Learned striatal age-decay rate versus published estimates.} (a) Horizontal range plots for Fearnley \& Lees 1991~\cite{fearnleyLees1991} (whole substantia nigra, clinicopathological), this paper's learned $k_{\mathrm{age}}$ across the twelve MLP-grid configurations, and Dzialas 2025~\cite{dzialas2025dat} (putamen-specific, DaT-SPECT longitudinal). Dots mark median values; bars span published or measured ranges. (b) Histogram of learned $k_{\mathrm{age}}$ across the twelve grid configs with the median indicated.}
\label{fig:p11-kage}
\end{figure}

\subsection{Subgroup-stratified performance}

To test whether the hybrid benefit is uniform across patient strata, we stratified the 65-patient test set by sex, age at baseline, and genetic risk. The test split contained 27 male and 38 female patients. Age groups split into 22 patients younger than $60$, 33 aged $60$--$70$, and 10 at $\geq 70$ years. We computed per-subgroup MAE for both the hybrid and the pure-mechanistic-fair baseline and the paired $\Delta$ with $95\%$ percentile bootstrap CI (Fig.~\ref{fig:p11-subgroup}).

The hybrid advantage is directionally consistent across all strata but varies in effect size. For males, the hybrid reduced MAE from $0.251$ to $0.168$ (paired $\Delta = -0.083$, $95\%$ CI $[-0.125, -0.045]$, excluding zero). For females, it reduced MAE from $0.172$ to $0.149$ ($\Delta = -0.024$, $95\%$ CI $[-0.077, +0.032]$, crossing zero). For patients younger than $60$, the hybrid reduced MAE from $0.232$ to $0.150$ ($\Delta = -0.081$, $95\%$ CI $[-0.145, -0.013]$, excluding zero). The $60$--$70$ and $\geq 70$ strata showed smaller point estimates of improvement whose confidence intervals crossed zero. The test set contained no LRRK2 or GBA carriers; genetic-risk stratification is therefore not reportable for this split, and is deferred to external validation cohorts where genetic-stratum-specific progression effects have been documented~\cite{iwaki2019gba}.

\begin{figure}[htbp]
\centering
\includegraphics[width=0.98\textwidth]{fig6_subgroup_mae.pdf}
\caption{\textbf{Subgroup-stratified test MAE: hybrid vs.\ pure-mechanistic-fair baseline.} Four panels stratify the 65-patient test set by sex, age group, LRRK2 carrier status, and GBA carrier status. Bars show per-subgroup MAE for the hybrid (blue) and pure-mechanistic-fair (grey) baselines with per-subgroup $n$ below each bar pair. The dotted reference line marks the cohort-level pure-mechanistic-fair MAE of $0.205$. The test set contained zero LRRK2 or GBA carriers; genetic stratification is not reportable on this split.}
\label{fig:p11-subgroup}
\end{figure}

Subgroup performance is broadly consistent across sex and age with no evidence of performance gaps that would flag equity concerns. The hybrid improvement is largest in precisely the strata where the pure-mechanistic baseline is weakest (males and younger patients). This pattern is consistent with the mechanism that the neural residual captures patient-specific heterogeneity that a single population decay rate cannot.

\subsection{Sample-size learning curve}

To assess whether performance had plateaued at cohort scale, we re-trained the best-cell hybrid ($\lambda_{\mathrm{phys}} = \lambda_{\mathrm{mono}} = 0.1$) at reduced training-set sizes of $n = 100$ and $n = 200$ patients, holding the $64$-patient validation and $65$-patient test splits fixed (Fig.~\ref{fig:p11-learning-curve}). Test MAE was $0.156$ at $n = 100$, $0.157$ at $n = 200$, and $0.152$ at $n = 299$ (baseline). The curve is near-flat between $n = 100$ and $n = 200$ and decreases modestly to $n = 299$. All three values remain substantially below the pure-mechanistic-fair MAE of $0.205$.

\begin{figure}[htbp]
\centering
\includegraphics[width=0.92\textwidth]{fig7_learning_curve.pdf}
\caption{\textbf{Learning curve for the hybrid UDE.} Hybrid test MAE (SBR units) as a function of training-set size at fixed best regulariser weights ($\lambda_{\mathrm{phys}} = \lambda_{\mathrm{mono}} = 0.1$). Validation and test splits are held fixed at 64 and 65 patients. The dashed horizontal line marks the pure-mechanistic-fair MAE of $0.205$. Test MAE is near-flat between $n = 100$ and $n = 200$ with a modest additional improvement at $n = 299$, suggesting the cohort is operating near the knee of the learning curve.}
\label{fig:p11-learning-curve}
\end{figure}

The curve shape has two implications. First, the physics prior and monotonicity regulariser are doing most of the work: a $100$-patient training fold already captures $99\%$ of the benefit of the full $299$-patient fold. Second, this cohort is operating near the knee of the learning curve. Substantially larger cohorts, such as those available from multi-site longitudinal DaT-SPECT collections, may yield modest further gains but are unlikely to produce order-of-magnitude improvements at this architecture.

\section{Discussion}

The headline result is a fair-baseline reversal. In an earlier accelerated-demo~\cite{dupre2026dissertation} the pure-mechanistic model appeared to beat the hybrid because it was given an oracle anchor, forecasting from the penultimate observed scan one step ahead. With the fair comparison --- both models forecasting from the same baseline initial condition over the same horizon --- the hybrid beats the pure-mechanistic baseline \emph{with the Fearnley--Lees literature prior}, which the diagnostic analysis (\S Rate misspecification) shows is structurally under-specified for regional SBR. The hybrid's value is that it learns the rate rather than requiring it a priori. Under real 5-fold cross-validation on all 428 patients (\S\ref{sec:robustness}) the pooled paired improvement is $\Delta = -0.053$ MAE ($95\%$ CI $[-0.065, -0.041]$), with all five fold-wise confidence intervals individually excluding zero, and the result is robust across all twelve regulariser-weight configurations. This closes the postdoc-scope question raised by the dissertation shallow-hybrid null. Physics-informed hybridisation \emph{can} extract additional signal from the data, but only when the physics shapes the predictor's internal dynamics rather than sitting beside the predictor's inputs.

At a single-fold test size of $n = 65$, the hybrid's paired $95\%$ CI against the trivial constant-mean predictor ($\Delta = -0.034$, $[-0.073, +0.013]$) marginally includes zero. The 5-fold cross-validated pooled analysis on all 428 patients (\S\ref{sec:robustness}), however, reports a $\Delta$ vs pure-mechanistic of $-0.053$, $[-0.065, -0.041]$ --- tight and cleanly excluding zero. The apparent borderline result at $n = 65$ is therefore a power limitation of the single-fold test, not a limitation of the architecture: the pooled evidence is unambiguous.

The architectural ablation clarifies which components are load-bearing. State-responsiveness --- letting the residual network see the solver's current state --- recovers half of the GRU-to-MLP performance gap. This suggests the learned dynamics at any instant depend on the current state, not just on patient covariates. Sequence memory beyond the current state confers no additional benefit at this cohort size. The simpler MLP residual is statistically indistinguishable from the best-tuned regularised GRU. This has practical implications. The MLP residual has approximately $10\times$ fewer parameters and trains in approximately half the wall-clock time, with no loss of predictive quality. Cohort sizes substantially larger than 428 may reveal a benefit to sequence memory; at PPMI scale, they do not.

The counterintuitive dropout--capacity interaction is consistent with the small-sample regime. At fixed network width, dropout acts as a capacity regulariser; if the baseline capacity is already near-critical for the task, dropout pushes the network below the effective floor and harms performance. We observed this directly: dropout $= 0.2$ applied to the GRU at $h = 32$ hidden units degraded performance by $0.005$ MAE with $95\%$ CI excluding zero. Reducing hidden size to $h = 16$ and then adding dropout recovered benefit. The practical implication is that in modest-cohort PD modelling, capacity and regularisation must be tuned jointly, not independently.

When initialized near the Fearnley--Lees anchor ($k_\text{age}^{(0)} = 0.025$/yr), the hybrid's learned $k_\text{age}$ converges to $0.045 \pm 0.013$/yr across 5 cohort folds --- within the published PPMI putamen SBR decline band ($9$--$14\%$/yr, see Marek~\cite{marek2015ppmi} and Nandhagopal~\cite{nandhagopal2009brain}). However, this recovery is anchor-conditional (\S\ref{sec:prior_sensitivity}): at other initializations in the learning regime, the learned rate doubles the anchor; in the plateau regime it barely moves. The hybrid's biological defensibility therefore rests on the anchor choice being in the ballpark, not on a claim of independent rate discovery. Our phenomenological whole-striatum fit lies between caudate ($2$--$3\%$/yr) and putamen ($4$--$6\%$/yr) region-specific rates, consistent with the aggregated region we model; region-stratified future work should target sub-region rates.

The subgroup analysis suggests that the hybrid's advantage is directionally consistent across demographic strata but largest where the pure-mechanistic baseline is weakest. Sex-stratified and age-stratified $\Delta$s are negative in every cell but only exclude zero for males and for patients younger than $60$. This pattern is consistent with a mechanism in which the neural residual absorbs heterogeneity the fixed population decay rate cannot capture. It is not consistent with the hybrid exploiting a subgroup-specific confound. Three points merit explicit mention. First, the sex effect may partially reflect baseline SBR differences between males and females in PPMI; the fair-mechanistic baseline starts further from the true endpoint for males, providing more room for the hybrid to improve. Second, the age effect is consistent with the literature observation that younger-onset PD exhibits faster decline~\cite{simuni2018longitudinal}, which a population-fixed $k_{\mathrm{age}}$ cannot capture but a covariate-conditioned residual can. Third, the absence of LRRK2 or GBA carriers in the $65$-patient test split is a genuine limitation of the random-split design. Stratified splitting, with carriers distributed across splits in fixed proportions, would address this in external-validation work.

The learning-curve result reframes the value of larger cohorts for this architecture. A $100$-patient training fold captures approximately $99\%$ of the benefit of the full $299$-patient fold at these regulariser settings. The mechanistic prior and monotonicity regulariser are therefore doing most of the work, and additional cohort scale is unlikely to yield order-of-magnitude improvements at this architecture. This is a positive finding for translational deployment: the hybrid framework does not require multi-thousand-patient training sets to operate. Multi-site DaT-SPECT cohorts would still be valuable for external validation and for testing architecturally richer variants.

Three interpretive caveats follow. First, our learning curve is traced at the best-cell regulariser setting only; other grid cells with weaker monotonicity priors might plateau differently, and a more exhaustive sweep at each training size is a natural extension. Second, the pure-mechanistic baseline does not learn and therefore has a training-independent floor of $0.205$. A data-hungrier hybrid would only appear disadvantaged at small $n$ if compared against the pure neural ODE rather than the mechanistic baseline. Third, because the validation split is held fixed across training sizes, the modest $n = 200$ bump above $n = 100$ reflects stochastic variation in which specific training patients are selected, rather than a real performance regression. Replicating the curve across multiple random seeds would tighten this interpretation, which we leave to future work.

Five limitations temper these claims. First, the analysis is single-cohort. External validation on longitudinal PD DaT-SPECT data is required before clinical deployment. Candidate cohorts include DeNoPa, SURE-PD3, and ICEBERG, identified as highest-priority external-validation targets in our companion NASEM-audit work~\cite{dupre2026paper10}. Existing PD progression analyses have generally relied on cross-cohort replication to establish robustness~\cite{severson2021progression, oxtoby2021ebm}. Second, the cohort size of 428 patients is modest. The near-flat learning curve suggests that cohort scale alone is unlikely to close the residual gap to a hypothetical biological ceiling. We cannot rule out that sequence-memory architectures would dominate at $n \gg 1000$. Third, we report only point-estimate MAE; uncertainty quantification on individual-patient ODE solutions is methodologically difficult and is not attempted here. Fourth, we predict the whole-putamen mean SBR; region-stratified six-ROI extension is a natural next step because the Dzialas analysis suggests region-specific rates differ substantially. Fifth, the three new analyses introduced here all carry their own caveats. The subgroup analysis is exploratory and underpowered for rare genetic carriers (zero in the random test split). The learning curve is traced at a single regulariser point and may differ at other grid cells. All results are internal-cohort. These three expansions should be considered preliminary until external validation is available.

Three concrete future directions follow from these limitations. First, external validation on longitudinal PD DaT-SPECT cohorts would establish cross-cohort generalisation. Second, uncertainty propagation through the ODE solver, via ensemble-based Monte Carlo or conformal post-hoc calibration, would produce clinically reportable prediction intervals at the individual-patient level. Third, a region-stratified six-ROI extension would apply the UDE separately to caudate and putamen substructures with a shared patient-covariate residual. This would test whether the architectural principle transfers beyond the whole-putamen summary. It would also enable direct comparison against Dzialas 2025~\cite{dzialas2025dat} at matched regional granularity.

\section{Methods}

\subsection{Data and split}

We used the PPMI longitudinal feature set~\cite{dupre2026paper3} from the Parkinson's Progression Markers Initiative~\cite{marek2018ppmi, brumm2023milestone}. After filtering to PD and prodromal patients with three or more putamen DaT-SPECT scans and complete baseline covariate data, 428 patients remained. Splits were patient-level random at 70/15/15 with seed 42, yielding 299/64/65 train/validation/test patients. No patient appeared in more than one split. The eleven baseline covariates that entered the residual network were age at visit, sex, MDS-UPDRS~\cite{goetz2008mdsupdrs} Part I total, MDS-UPDRS Part III total, Epworth Sleepiness Scale total, REM-Sleep-Behavior-Disorder Screening Questionnaire total, SCOPA-AUT total, UPSIT total, MoCA total, LRRK2 carrier status, and GBA carrier status. The target was the putamen mean SBR at each patient's held-out last scan.

\subsection{Models}

Three architectures were evaluated. The pure-mechanistic-fair baseline solves $dS/dt = -k_{\mathrm{age}} S$ with fixed $k_{\mathrm{age}} = 0.025$/yr over the full horizon from each patient's baseline observed SBR; this is the fair comparison because it starts from the same initial condition as the hybrid. The pure neural ODE replaces the dynamics with $dS/dt = \mathrm{NN}(S, x; \theta)$ using a multilayer-perceptron residual of width 32 with a single hidden layer and SiLU activation. The physics-informed hybrid UDE~\cite{rackauckas2020universal, chen2018neuralode} uses $dS/dt = -k_{\mathrm{age}} S + \mathrm{NN}_{\mathrm{res}}(S, x; \theta)$, where $k_{\mathrm{age}}$ is a learned positive scalar anchored at the Fearnley--Lees prior $0.025$/yr~\cite{fearnleyLees1991} via a soft Gaussian prior, and the neural residual is initialised small (output-layer zero bias, gain $0.1$ on weights) so that the physics term dominates at initialisation. All ODEs are integrated forward with a fixed-step fourth-order Runge--Kutta scheme at $\Delta t = 0.25$ years using torchdiffeq~\cite{chen2021torchdiffeq} on top of PyTorch~\cite{paszke2019pytorch}.

\subsection{Loss and regularisation}

The total loss is given by Eq.~(1):
\begin{equation}
\mathcal{L}(\theta) = \mathrm{MSE}(S_{\mathrm{obs}}, \hat{S}) + \lambda_{\mathrm{phys}} \cdot \|\mathrm{NN}_{\mathrm{res}}\|_2^2 + \lambda_{\mathrm{mono}} \cdot \mathbb{E}_t\!\left[\mathrm{ReLU}(dS/dt)\right],
\label{eq:loss}
\end{equation}
where the physics penalty $\lambda_{\mathrm{phys}}$ suppresses the neural residual towards zero (driving the dynamics towards pure mechanistic), and the monotonicity penalty $\lambda_{\mathrm{mono}}$ discourages the solver trajectory from exhibiting positive gradients (consistent with SBR typically declining in the PD direction over the relevant timescale). We swept $\lambda_{\mathrm{phys}} \in \{0, 0.001, 0.01, 0.1\}$ and $\lambda_{\mathrm{mono}} \in \{0, 0.01, 0.1\}$ for a full $4 \times 3 = 12$ configuration grid. The monotonicity regulariser follows the physiological-UDE formulation of de Rooij and colleagues~\cite{deRooij2025physiology}.

\subsection{GRU residual variants}

For the architectural ablation we replaced the MLP residual with a gated recurrent unit~\cite{cho2014gru} in four variants at fixed $\lambda_{\mathrm{phys}} = 0.01$, $\lambda_{\mathrm{mono}} = 0.01$. The vanilla GRU (v1) consumes only patient covariates $x$ and emits the residual; it has no access to the solver's current state $S$ during the forward pass. The state-aware GRU (v2) appends the solver's current state $S$ as a synthetic final sequence step to the GRU input sequence at each ODE evaluation. The small state-aware variant reduces the GRU hidden size from 32 to 16. The small-plus-dropout variant additionally applies dropout with probability $0.2$ to the GRU output before the residual projection.

\subsection{Training}

All models were trained with Adam~\cite{kingma2014adam} at learning rate $10^{-3}$ (batch size 16, grouped by patient) for up to 100 epochs with early stopping on validation MAE at patience 10. Training used a single CPU without distributed parallelism; each configuration took approximately 15--25 minutes wall-clock to converge. The learning-curve variants at $n = 100$ and $n = 200$ were trained with the same hyperparameters, differing only in a deterministic sub-selection of the 299-patient training pool controlled by the \texttt{--train-cap} flag.

\subsection{ODE-solver sensitivity}

To verify that the reported test MAE is not sensitive to ODE-solver discretisation, we ran the winning configuration with four different solver configurations: dopri5 with default tolerances ($\text{rtol}=10^{-3}$, $\text{atol}=10^{-4}$); dopri5 with tight tolerances ($\text{rtol}=10^{-5}$, $\text{atol}=10^{-6}$); RK4 at a fixed step size of $0.1$ years; and dopri8 adaptive. All four solvers produced test MAE within $0.0001$ of one another ($0.1522$--$0.1523$; Table~\ref{tab:p11_solver_sensitivity}), with per-patient maximum absolute error difference of $5.7 \times 10^{-4}$ SBR units (0.4\% relative to the cohort mean error) and Wilcoxon paired-test $p$-values between $0.22$ and $0.36$. The ODE is numerically well-converged at dopri5 default tolerances for PPMI DaT-SPECT trajectory scales.

\begin{table}[ht]
\centering
\caption{Solver sensitivity. Winning configuration retrained under four solver regimes. Test MAE invariant to four decimals.}
\label{tab:p11_solver_sensitivity}
\begin{tabular}{lc}
\toprule
Solver & Test MAE \\
\midrule
dopri5 default ($\text{rtol}=10^{-3}$, $\text{atol}=10^{-4}$) & $0.1523$ \\
dopri5 tight ($\text{rtol}=10^{-5}$, $\text{atol}=10^{-6}$) & $0.1522$ \\
RK4 fixed-step ($dt = 0.1$ yr) & $0.1522$ \\
dopri8 adaptive & $0.1522$ \\
\bottomrule
\end{tabular}
\end{table}

\subsection{Evaluation}

Primary metric was held-out last-scan MAE. Paired bootstrap with 1000 resamples (with replacement, sampling test patients) produced the $95\%$ confidence interval on $\Delta$(hybrid $-$ pure-mechanistic-fair); we report the $2.5$ and $97.5$ percentiles of the bootstrap distribution. Paired bootstrap resamples patients (not sentences or configurations), preserving patient-level pairing of per-model error terms. For the subgroup analysis, the same paired-bootstrap procedure was applied within each stratum. The constant-mean baseline is computed as a single constant equal to the training-set mean last-scan SBR. For the primary robustness evaluation we implemented real 5-fold cross-validation: the 428-patient cohort was deterministically partitioned into five non-overlapping test folds (master seed $=42$), the winning configuration was retrained from scratch on each fold's train+val set, and the 428 patient-level predictions were pooled across folds for the primary $\Delta$ bootstrap. Per-fold summary statistics (MAE mean $\pm$ SD, range, t-distribution CI on the mean) are reported alongside the pooled estimate. The Option-B holdout-discipline rerun (seed 2026) is retained as a secondary single-split confirmation. The winning configuration is not re-tuned on any holdout. Full code, per-patient predictions, and grid-level summaries are archived in PostgreSQL tables \texttt{mechanistic.paper11\_sciml\_summary} (including the \texttt{kfold5\_lp0.1\_lm0.1\_mlp\_fold\{0..4\}} entries) and \texttt{mechanistic.paper11\_sciml\_results}.

\section*{Data Availability}

PPMI data are available to qualified researchers through the Michael J. Fox Foundation at \url{https://www.ppmi-info.org/access-data-specimens/download-data}, with a required data use agreement. Details of cohort selection, staging, and feature assembly follow the Parkinson's Progression Markers Initiative design~\cite{marek2018ppmi, brumm2023milestone}. Derived cohort definitions, model predictions, and summary statistics specific to this paper are archived in the companion dissertation~\cite{dupre2026dissertation} local PostgreSQL database (tables \texttt{mechanistic.paper11\_sciml\_summary} and \texttt{mechanistic.paper11\_sciml\_results}). A reproducibility snapshot is available from the corresponding author upon request pending defense completion.

\section*{Code Availability}

Producer (\texttt{scripts/paper11\_demo/hybrid\_sciml\_full\_cohort.py}, 998 lines) and database loader (\texttt{scripts/load\_paper11\_to\_pg.py}) are version-controlled at \url{https://github.com/bddupre92/PD\_PHD}, branch \texttt{feat/ch9-6-multichannel}. Public release is scheduled upon defense completion. Independent reproduction using PPMI data requires a PPMI data use agreement and Python 3.12+ with PyTorch 2.8+ and torchdiffeq.

\section*{Author Contributions}

B.D.D. conceived and designed the study (Conceptualisation), implemented the physics-informed neural ODE and baselines (Software, Methodology), ran the full-cohort scale-up and architectural ablations (Formal analysis), generated the publication figures (Visualisation), and wrote the manuscript (Writing -- original draft, Writing -- review \& editing).

\section*{Competing Interests}

The author declares no competing interests.

\section*{Ethics Declarations}

This work analyses de-identified data from the Parkinson's Progression Markers Initiative (PPMI), which is approved by the institutional review boards (IRB) of all participating sites~\cite{marek2018ppmi}. No additional IRB review was required for this secondary analysis of de-identified data.
